%----------------------------------------------------------------------------------------
% Tailored CV — Doctoral Researcher, AI-Driven Protein Design
% CIIS / Chair Prof. Detlef Schoder, Universität zu Köln — Ref. Wiss2604-xx
%----------------------------------------------------------------------------------------

\documentclass[10pt,a4paper,sans]{moderncv}

\usepackage[utf8]{inputenc}
\usepackage[T1]{fontenc}

\moderncvstyle{classic}
\moderncvcolor{blue}

\usepackage[scale=0.78]{geometry}
\usepackage{ragged2e}

%----------------------------------------------------------------------------------------

\firstname{Kundai Farai}
\familyname{Sachikonye}

\address{Auf der Lichtung 19}{82178 Puchheim, Germany}
\mobile{(+49) 1626386344}
\email{kundai.sachikonye@bitspark.com}
\homepage{github.com/fullscreen-triangle}

%----------------------------------------------------------------------------------------

\begin{document}

\makecvtitle

%----------------------------------------------------------------------------------------
%	EDUCATION
%----------------------------------------------------------------------------------------

\section{Education}

\cventry{2019--2021}{PhD candidate, Computational Lipidomics}{Technische Universität München / Bayerisches Forschungsinstitut}{}{}{
  Molecular profiling at HPC scale, multi-omics statistical modelling,
  physics-based feature engineering. Departed to pursue independent research.
}
\cventry{2018--2019}{MSc Computational Biology}{Universität Tübingen}{}{}{
  Protein structure and dynamics, molecular modelling, computational biophysics,
  ML for biological sequences.
}
\cventry{2014--2017}{MSc Pharmaceutical Biotechnology}{HAW Hamburg}{}{}{}
\cventry{2011--2014}{BSc Biochemistry and Cell Biology}{Jacobs University Bremen}{}{}{}

%----------------------------------------------------------------------------------------
%	RELEVANT MANUSCRIPTS
%----------------------------------------------------------------------------------------

\section{Relevant Manuscripts}

\cvitem{2025}{
  \textbf{On the Thermodynamic Consequences of Categorical Completion on Biological Systems.}
  MHC-I/II binding affinity from spectral complementarity: AUC 0.81 vs.\ NetMHCpan 0.68
  ($r=0.73$, $p<10^{-12}$, $n=2{,}847$ IEDB peptides). Zero-parameter PPI binding
  prediction. \textit{Manuscript.}
}

\cvitem{2025}{
  \textbf{Backward Training Principle.}
  Formal proof: training at the pure-intent regime (minimum binding energy,
  maximum on-target specificity) propagates $\epsilon$-competence to all
  constrained design sub-tasks by feasibility inclusion.
  Forward Training Collapse Theorem (positive VC failure gap for conventional
  direction); Sample Complexity Theorem ($(N{+}1)\times$ savings).
  \textit{Manuscript.}
}

%----------------------------------------------------------------------------------------
%	PROTEIN DESIGN AND PPI TOOLS
%----------------------------------------------------------------------------------------

\section{Protein Design and PPI Tools}

\cvitem{crown-prince}{
  \textbf{Zero-Parameter Physics-Informed Binding Affinity Scoring.}
  $K_d = \exp(-\Delta M \cdot \ln b)$ from bounded phase space partition geometry;
  Hill equation as exact limiting case; PD/PK/ADME trajectories from partition structure.
  Directly integrable as physics-informed scoring function in ML protein design pipelines.
  Validated: 39/39 NIST binding compounds, zero free parameters; 50\,MB vs DrugBank 12\,GB.
  Cell Spectral Hologram: tumour cell $\to$ spectral embedding $\to$ therapeutic
  perturbation vector; AUC$=1.00$.
  Live: \href{https://crown-prince-four-courners.vercel.app/tools/spectrum}{crown-prince-four-courners.vercel.app}.
  \href{https://github.com/fullscreen-triangle/crown-prince}{github.com/fullscreen-triangle/crown-prince}
}

\cvitem{syndrome}{
  \textbf{Spectral Physics-Based PPI Prediction and Lead Optimisation.}
  3-channel DFT (hydropathy, volume, charge) $\to$ 36D cosine similarity $\rho$;
  $\mathcal{O}(cL\log L)$; no training data required.
  \textbf{PPI binding:} $\rho\geq0.72$; AUC 0.81 vs.\ NetMHCpan 0.68 on 2,847 IEDB peptides.
  \textbf{Lead optimisation:} escape predictor screens all single-aa substitutions for
  variants maximising target $\rho$ and minimising off-target $\rho$;
  score $E(v) = -\max\rho(v,\mathrm{off}) + w\cdot\rho(v,\mathrm{target})$.
  Physics-based feature layer for hybrid scoring frameworks.
  Live: \href{https://syndrome-seven.vercel.app/tools}{syndrome-seven.vercel.app/tools}.
  \href{https://github.com/fullscreen-triangle/syndrome}{github.com/fullscreen-triangle/syndrome}
}

\cvitem{borgia}{
  \textbf{Zero-Parameter Molecular Feature Generation.}
  Electronic transitions, absorption/emission spectra, vibrational modes, and
  partition coordinates from bounded phase space geometry --- zero empirical parameters.
  Triple Equivalence: oscillatory $=$ categorical $=$ partition; GPU shader $=$
  physical spectrometer. Generates physics-informed input representations for
  protein design ML models that generalise by construction to out-of-distribution sequences.
  Validated: NIST 9 elements + 6 molecules, self-consistency $<$1.63\%.
  Live: \href{https://honjo-masamune.vercel.app/tools}{honjo-masamune.vercel.app/tools}.
  \href{https://github.com/fullscreen-triangle/borgia}{github.com/fullscreen-triangle/borgia}
}

\cvitem{shakespear}{
  \textbf{Protein Conformation, Trajectory, and Dynamics Analysis.}
  S-entropy coordinates $(S_k, S_t, S_e)\in[0,1]^3$ for protein conformational states;
  trajectory completion operator $T=P\circ B\circ K$ (Kirchhoff propagation,
  backward Viterbi, thermodynamic projection); Banach contraction, unique fixed-point.
  Lightweight fold evaluation and transition-state characterisation without
  full MD or AlphaFold inference.
  Live: \href{https://shakespear-nine.vercel.app/instrument}{shakespear-nine.vercel.app/instrument}.
  \href{https://github.com/fullscreen-triangle/shakespear}{github.com/fullscreen-triangle/shakespear}
}

\cvitem{purpose + homo-veloce}{
  \textbf{Backward Training Principle for Protein Design Optimisation.}
  Formal proof: training at the pure-intent binder regime (minimum $K_d$, maximum
  on-target specificity, no extrinsic constraints) propagates $\epsilon$-competence
  to all constrained design sub-tasks (solubility, stability, off-target avoidance)
  by feasibility inclusion.
  homo-veloce distils physics-based solvers into lightweight neural surrogates
  (1D CNNs, LSTMs, Transformers) --- the same distillation pattern as training
  ML design models on physics-informed features.
  \href{https://github.com/fullscreen-triangle/purpose}{github.com/fullscreen-triangle/purpose}
}

%----------------------------------------------------------------------------------------
%	AGENTIC AI SYSTEMS
%----------------------------------------------------------------------------------------

\section{Agentic AI Systems}

\cvitem{four-sided-triangle}{
  \textbf{8-Stage Metacognitive Pipeline.}
  Adaptive routing between LLM reasoning and physics/mathematical solvers;
  Rust backend (10--50$\times$); FastAPI; Docker/Kubernetes; GQIC resource allocation.
  Production template for protein design agent: sequence proposal $\to$ physics
  scoring $\to$ structural validation $\to$ feedback loop.
  \href{https://github.com/fullscreen-triangle/four-sided-triangle}{github.com/fullscreen-triangle/four-sided-triangle}
}

\cvitem{bloodhound + zangalewa}{
  \textbf{Distributed Agentic VM and Routing Architecture.}
  Bloodhound: Triangle language (navigate/slice/complete); St-Hurbert engine;
  Common-Cell Convergence Theorem; Intelligence Index $O(1)$.
  Zangalewa: Minimum Sufficient Interceptor; $O(\log_k N)$ specialist cascade;
  Sufficiency Theorem (parameter count independent of substrate size).
  \href{https://github.com/fullscreen-triangle/bloodhound}{github.com/fullscreen-triangle/bloodhound}
}

%----------------------------------------------------------------------------------------
%	EXPERIENCE
%----------------------------------------------------------------------------------------

\section{Experience}

\cventry{2021--present}{Independent AI \& Computational Biology Developer}{Bitspark GmbH}{}{}{
  Physics-informed scoring for protein-ligand and protein-protein binding;
  protein design ML pipelines; agentic system architecture; all frameworks
  publicly deployed and validated.
}

\cventry{2019--2021}{Computational Lipidomics}{TUM / Bayerisches Forschungsinstitut}{Munich}{}{
  MS-based molecular profiling, multi-omics integration, distributed clinical pipelines.
  Python, Java.
}

\cventry{2019}{Process Optimisation}{Fraunhofer Institut IGB}{Stuttgart}{}{
  ML-based optimisation for biological systems. Python, C.
}

%----------------------------------------------------------------------------------------
%	TECHNICAL SKILLS
%----------------------------------------------------------------------------------------

\section{Technical Skills}

\cvitem{Protein Design / Structural Biology}{
  Physics-informed scoring (zero-parameter $K_d$, spectral $\rho$ PPI binding);
  protein conformation analysis; S-entropy trajectory completion;
  molecular feature generation (electronic transitions, vibrational modes);
  AlphaFold-compatible structural feature integration; RDKit cheminformatics
}

\cvitem{ML / DNN}{
  PyTorch; 1D CNNs, LSTMs, Transformers, GNNs; knowledge distillation
  from physics solvers; Bayesian networks; variational Bayes;
  Gaussian processes; surrogate model training; RAG pipelines
}

\cvitem{Agentic AI}{
  Multi-stage metacognitive pipelines; specialist cascade routing $O(\log_k N)$;
  COMPILE/EXECUTE two-phase agents; Intelligence Index agent auditing;
  LLM orchestration (Claude, OpenAI, Ollama); tool use and function calling
}

\cvitem{Languages}{
  Python (primary), Rust, Java, TypeScript, C, Bash, SQL, R
}

\cvitem{Infrastructure}{
  Ray, Dask; Docker, Kubernetes, FastAPI; WebGPU compute shaders;
  memory-mapped I/O; HDF5; Git, CI/CD
}

%----------------------------------------------------------------------------------------
%	LANGUAGES
%----------------------------------------------------------------------------------------

\section{Languages}

\cvitemwithcomment{English}{Native / Fluent}{}
\cvitemwithcomment{German}{Conversational}{Living in Germany since 2011}

%----------------------------------------------------------------------------------------

\renewcommand{\listitemsymbol}{-~}

\end{document}
